\documentclass[11pt]{article}
\usepackage[margin=1in]{geometry}
\usepackage[T1]{fontenc}
\usepackage{lmodern,amsmath,amssymb,booktabs,graphicx}
\usepackage[colorlinks=true,linkcolor=blue,urlcolor=blue]{hyperref}
\setlength{\parindent}{0pt}\setlength{\parskip}{5pt}
\setlength{\emergencystretch}{2em}
\title{RoCE-K: confirmation, evaluated design removal, and model boundaries}
\author{}\date{4 October 2026}
\begin{document}\maketitle

\textbf{Summary.} Independent seeds confirm the previously observed length
advantage at $p=4,K=2048$ with constant target effect. Actual design removal
can recover borrowing when the retained designs are stable, but eligibility
alone is insufficient. A separate model-class comparison shows that the
current bounded-invalid certificate can be substantially conservative. A new
conditional worst-case lower bound explains why a uniform fourth-root gain
over the entire unrestricted-invalid class nevertheless cannot be promised.

\section{Accounting and reproducibility}

The new panels contain 4,400 outcome setting-specific repetitions: four
independent confirmation settings, ten paired design-policy settings, and
eight paired model-boundary settings, each with 200 repetitions. No failed
repeat is replaced. Multiple methods share patient data within a repetition;
nonlinearity amplitudes also share seeds within each boundary source count.
These are not 4,400 mutually independent datasets. The earlier 200-repeat
discovery results are retained and analyzed separately.

Confirmation uses seed offset 30190000 and the unchanged version-7 complete-data
runner. Design removal and boundary panels use offsets 31190000 and 32190000.
Conditional intervals, population intervals, design failures, retained indices,
adjusted valid counts, source noise and dispersion contributions are saved.
Every stored summary is checked against interval endpoints. The new selection
function accepts only covariates and treatment assignments; outcome evaluation
occurs after the retained set is fixed.

The research tests now total 62 passing checks, including manual-subset parity,
outcome-independent selection, exhaustive valid-set count inequalities,
projection-budget recomputation, exact mixture identities, and agreement of
grouped sufficient statistics with patient-level leave-out calculations.
Production RoCE/ENAR, previous results and the version-7 manuscript are preserved.

\section{Independent confirmation and paired Monte Carlo uncertainty}

The four confirmation settings retain 1,000 patients per site, bounded
continuous covariates, four signal coordinates, source covariate tilt .4 and
population TATE .10. The declared valid fraction stays at $3/4$, including
the all-valid control. Weak invalid treated-arm shifts are
$2\sqrt{.5/1000}\,K^{-1/4}$. No true constant-effect information enters the
target coefficient certificate. All methods use exactly the same draws
within a setting. The difference column is source minus ordinary target-normal
length, and its standard error uses paired differences.

\begin{center}\small
\resizebox{\linewidth}{!}{\input{experiments/20261004c/confirmation_table}}
\end{center}

For $p=4,K=2048,h=0$ with weak bias, the old discovery cohort had paired mean
length difference $-.010898$ with Monte Carlo standard error $.000593$.
The independent confirmation gives $-.010421$ with standard error $.000607$,
or an 8.49\% reduction relative to ordinary normal target intervals. The
source and target-normal observed coverages are $1.000$ and $.930$; their
95\% exact binomial intervals are $[.9817,1]$ and $[.8853,.9612]$.
Thus the length finding reproduces, while 200 draws do not establish that
the normal comparator's true coverage differs from $.95$.

The all-valid control has a similar 8.44\% length reduction and coverages
$1.000$ versus $.955$. The $p=10$ and heterogeneous-target controls do not
beat ordinary normal target length. All four source procedures remain shorter
than the equally protected target-only comparator. This is a setting-specific
efficiency result, not a general optimality claim.

\section{Evaluated design-only removal}

Both policies use the same rank, condition-number and leverage guards. A site
failing either arm is removed in its entirety. With $m$ removed sites, inference
uses $K'=K-m$ and $g_a'=\max(0,3K/4-m)$. Every source normalization and
dispersion/approximation constant is recalculated; the invalid-count allowance
is $K'-g_a'$, not $K'/4$. If either adjusted valid count is zero, inference
uses the protected target band. All designs are selected without outcomes.

The first stress family has $n_t=1000$, source sizes 120 or 200, $p=45$,
and $\Pr(A=1\mid X)=\operatorname{expit}\{2.5(X_1+X_2)\}$ at sources.
Target treatment is randomized with probability $.5$. Outcome means remain
bounded and linear, with target heterogeneity $.15$. This deliberately severe
design stress is separate from the 1,000-patients-per-site confirmation panel.

\begin{center}\small
\resizebox{\linewidth}{!}{\input{experiments/20261004c/removal_table}}
\end{center}

The table uses weak-bias settings. At $n_s=120,p=45,K=64$, the adjusted-count
policy is eligible in every repetition, yet its mean width equals the original
policy's $.3875$. The source interval remains too wide to improve on the target
intersection. Raising $n_s$ to 200 removes observed design failures but still
does not produce useful borrowing in this small-$K$, poorly overlapping design.
This directly distinguishes count eligibility from precision.

The second family retains 1,000 patients per site and $p=4,K=256$, with
randomized treatment and target heterogeneity $.15$. In its prespecified
design-failure setting, the final eighth of sources has $X_4=X_3$; these
sources are structurally valid but cannot balance the full target feature
map. The failure mechanism is fixed before outcomes. Removing the 32 failed
sites leaves $K'=224$ and conservative $g_a'=160$.

Here length decreases from $.3092$ to $.1617$, with paired difference
$-.14751$ and Monte Carlo standard error $.000323$. The improved protected
target-only comparator has length $.2825$. Bounded-invalid inference on the
retained set gives $.1756$, preserving a substantial gain while protecting a
broader class. The removal methods have observed conditional coverage $.995$
and population coverage $1.000$. With no deliberately singular sites, the two
policies are identical, as required. The all-valid stress controls, all design
diagnostics and the normal target comparison are included in the saved tables.

\section{Which part of the broader-model floor is structural?}

The first calculation is algorithm-specific. With $N$ observations per arm,
uniform balance weights and declared invalid fraction $\varepsilon$, the
current bounded-invalid certificate has limiting source bias radius
\[
 \frac{\varepsilon}{\sqrt{(1-\varepsilon)(N-1)}}
\]
even if all sources actually have identical valid constant means. This is a
property of the certificate, not an all-method adaptive lower bound.

The paired experiment uses the earlier three-point fixed design and Bernoulli
outcomes, $N=500$ per arm. The two prespecified bases are $(1,x)$ with
bounded-invalid protection and $(1,x,x^2)$ with exact leave-out correction.
The latter spans every mean function on this finite support and requires no
knowledge of the nonlinearity amplitude. It still retains unknown weak
invalid-source mean shifts and the same valid-count assumption.

\begin{figure}[ht]\centering
\includegraphics[width=\linewidth]{experiments/20261004c/model_boundary.pdf}
\caption{Paired conditional source bands on three-point support. Source count
increases while each arm retains 500 observations. No population enlargement
is added to these comparisons.}\end{figure}

At $K=262144$ and amplitude $.4$, conditional interval length is $.02575$
under bounded-invalid protection with the linear basis, versus $.00547$ with
the saturated basis. Conditional coverages are $.995$ and $1.000$; source-band
coverages are both $1.000$. At $K=32768$, saturated conditional/source coverage
is $.995$ for both amplitudes. These results show that declared design/model
structure can remove much of the certificate cost. Repeated finite support is
additional structure; it cannot be silently assumed for unrestricted continuous
invalid means.

The second calculation is a distinct worst-case lower bound in the main text.
On distinct fixed covariates, arbitrary bounded invalid means can realize a
prior over independent degenerate Bernoulli patient laws. Two different valid
centers can then have identical contaminated source-data mixtures. Conditioning
these priors to contain at least $g$ valid sites costs an explicit binomial tail
$q_K$. Including the trusted target yields the interval-length bound
\[
 \sup_P E_P|I|\ge
 \Delta\,[1-2\alpha-t_t-2q_K]_+.
\]
For $N_s=N_t=500$, contamination probability $.2$, guaranteed valid fraction
$.75$, $K=2048$, and $\Delta=.1/\sqrt{500}$, exact calculation gives
$t_t=.0796502$, $q_K=1.60\times10^{-8}$, and lower bound $.0036687$.
The mixture identity is verified to floating-point precision.

This establishes a conditional worst-case obstruction of order $n_s^{-1/2}$
when source and target sizes are comparable. It does not prove the current
certificate optimal, does not concern an all-valid adaptive distribution,
and is not yet an i.i.d.-random-design population lower bound. It also does
not replace a lower bound specifically restricted to weak-shift configurations.

\section{What remains}

The main text now gives one algorithm, an explicit error ledger and a single
coverage theorem with model-specific length conditions. The next theoretical
question is honest adaptation inside the broader class: can a procedure attain
shorter intervals at favorable distributions while still protecting arbitrary
bounded invalid means? The new worst-case construction does not answer that.
For practical development, retained-design stability and target coefficient
protection remain the relevant precision costs. Any further source screening
must remain design-only unless its outcome selection uncertainty is separately
calibrated.

\end{document}
